% Generated draft: requires mathematical and source review.
\documentclass[11pt]{article}
\usepackage[margin=25mm]{geometry}
\usepackage{fontspec}
\setmainfont{Latin Modern Roman}
\usepackage{amsmath,amssymb,mathtools,unicode-math}
\setmathfont{Latin Modern Math}
\usepackage{enumitem,needspace,xurl,hyperref}
\setlength{\parindent}{0pt}
\setlength{\parskip}{5pt}
\setlength{\emergencystretch}{3em}
\newcommand{\N}{\mathbb{N}}
\newcommand{\Z}{\mathbb{Z}}
\newcommand{\Q}{\mathbb{Q}}
\newcommand{\R}{\mathbb{R}}
\newcommand{\C}{\mathbb{C}}
\newcommand{\F}{\mathbb{F}}
\newcommand{\A}{\mathbb{A}}
\newcommand{\PP}{\mathbb{P}}
\newcommand{\E}{\mathbb{E}}
\newcommand{\eps}{\varepsilon}
\newcommand{\dd}{\,\mathrm{d}}
\begin{document}
\section*{UnsolvedMath ID: 30004687}
\textbf{Source problem code:} catalogue entry
\subsection{Definitions and mathematical statement}
For a complex projective variety \(X\), let \(X^{\mathrm{an}}\) denote its associated complex analytic space. An entire curve in \(X\) is a nonconstant holomorphic map \(f:\C\to X^{\mathrm{an}}\). Define the entire-curve locus and its Zariski closure by
\[
E(X)=\bigcup_{\substack{f:\C\to X^{\mathrm{an}}\ \mathrm{holomorphic}\\ f\ \mathrm{nonconstant}}}f(\C),
\qquad
Z_{\mathrm{ent}}(X)=\overline{E(X)}^{\mathrm{Zar}}.
\]
Here the union is empty if there are no entire curves. For a smooth irreducible complex projective variety of dimension \(n\), being of general type means that its canonical bundle \(K_X\) is big, equivalently that its Kodaira dimension satisfies \(\kappa(X)=n\). The supplied record does not specify smoothness, irreducibility, or a convention for general type in the singular case; these scope qualifications remain unresolved.

\textbf{Green--Griffiths--Lang question as supplied.} Is it true that, for every complex projective variety \(X\) of general type,
\[
\exists\,Y\subsetneq X\ \text{Zariski closed and algebraic}\quad
\forall\,f:\C\to X^{\mathrm{an}}\ \text{nonconstant and holomorphic},
\qquad f(\C)\subseteq Y?
\]
The same \(Y\) must work for every entire curve in \(X\). Interpreting the source's word ``subvariety'' as a Zariski-closed algebraic subset, allowing reducibility and the empty set, this is equivalent to
\[
Z_{\mathrm{ent}}(X)\ne X.
\]
Under this interpretation, \(Z_{\mathrm{ent}}(X)\) is the unique smallest Zariski-closed algebraic subset containing all entire-curve images. The source does not explicitly specify these conventions for ``subvariety.''

\textbf{Arithmetic qualification retained from the source.} The record additionally states that, over number fields, the smallest such \(Y\) is expected to coincide with the ``Mordell locus.'' No definition of that locus, field of definition, choice of complex embedding, or quantification over finite field extensions is supplied. Consequently, this arithmetic expectation is retained as an unresolved qualification, rather than formulated as a precise equality with an invented arithmetic locus.
\subsection{Short English statement}
The question asks whether all nonconstant entire curves in any complex projective variety of general type lie in one proper algebraic subset. The source also expects this smallest subset, in a number-field setting, to agree with a Mordell locus whose definition is not supplied. The singular-case scope and arithmetic comparison require clarification.
\subsection{Sources}
\begin{itemize}
\item Supplied local registry record, permanent UnsolvedMath ID \(30004687\), exact source problem code \texttt{catalogue entry}, in the Oberwolfach Reports Open Problems collection. Project registry: \texttt{data/unsolvedmath/problems.json}; supplementary registry: \texttt{catalog/problems.json}. No individual proposer is named; the supplied proposed year is 2021.
\item Supplied source citation: \emph{Geometric Methods of Complex Analysis (2022)}, Oberwolfach Reports, DOI 10.4171/owr/2021/25, pp.\ 1291--1345, \url{https://doi.org/10.4171/owr/2021/25}. The literature-source list also labels this reference ``Oberwolfach Report 25/2021.'' The supplied background describes Oberwolfach Reports as documenting mathematical workshops and their open-problem discussions.
\item Supplied literature reference, labeled ``Cadorel, Demailly and collaborators, hyperbolicity'': \url{https://doi.org/10.1007/s00222-023-01219-z}. No fuller bibliographic identification is supplied.
\item Supplied literature reference, labeled ``Demailly, recent results on GGL'': \url{https://arxiv.org/abs/1801.04765}.
\item Supplied literature assessment, checked \textbf{2026-08-22}: the registry labels the problem \texttt{partially\_solved} and states that the conjecture remains open for arbitrary varieties of general type, with major partial results for broad classes such as sufficiently positive generic hypersurfaces and varieties satisfying additional positivity or group-theoretic hypotheses. This is supplied dated curation derived from the OWR record, not an independently verified literature assessment; its status and cited claims require independent verification.
\item The registry labels the statement \texttt{corrected\_verified}, explaining that only formula delimiters, notation, grammar, and typography were repaired, with mathematical content and scope preserved. This supplied editorial label does not constitute independent mathematical verification of this draft.
\end{itemize}
\end{document}
